\documentclass[11pt]{article}
\usepackage[margin=1in]{geometry}
\usepackage{amsmath,amssymb,amsthm,booktabs}
\usepackage[colorlinks=true,urlcolor=blue,linkcolor=blue]{hyperref}
\newtheorem{theorem}{Theorem}
\newtheorem{lemma}{Lemma}
\newcommand{\R}{\mathbb R}
\title{Strict ratio log-concavity of Sun's \(R\)-sequence}
\author{Independent working note --- MathDB 331960}
\date{12 September 2026 (UTC)}
\begin{document}
\maketitle

\paragraph{Task 1: solve the published conjecture.}
Define
\[
 R_n=\sum_{k=0}^n\binom nk\binom{n+k}k\frac1{2k-1},
 \qquad r_n=\frac{R_{n+1}}{R_n}\quad(n\ge1).
\]
Conjecture 4.7 of Brian Yi Sun's published paper \cite{sun} asserts that
\(\{r_n\}_{n\ge4}\) is log-concave. We prove the strict assertion, with
an exact finite algebraic certificate for the induction.

\begin{theorem}\label{main}
For every integer \(n\ge5\),
\[
 r_n^2>r_{n-1}r_{n+1},
 \quad\text{equivalently}\quad
 R_{n+1}^3R_{n-1}-R_n^3R_{n+2}>0.
\]
Consequently the published Conjecture 4.7 is true, and its starting
index \(4\) cannot be decreased.
\end{theorem}

\paragraph{Indexing.}
The MathDB transcription defines its ratio as \(R_n/R_{n-1}\) while
retaining the starting index \(4\). In that notation the correct
starting index is \(5\). This note proves the original published
conjecture; it does not count the transcription error as a refutation.

\paragraph{A recurrence derived from the definition.}
For \(n\ge1\), the terms \(k=0,1\) give \(R_n\ge n(n+1)-1>0\).
Formal power series give
\[
 \sum_{n\ge0}R_nz^n
 =\frac1{1-z}\sum_{k\ge0}\frac{\binom{2k}k}{2k-1}
                  \left(\frac{z}{(1-z)^2}\right)^k
 =-\frac{\sqrt{1-6z+z^2}}{(1-z)^2}=:G(z).
\]
Here the middle identity follows from the binomial expansion
\(\sum_{k\ge0}\binom{2k}k\,u^k/(2k-1)=-\sqrt{1-4u}\).
Differentiating gives
\[
 (1-z)(1-6z+z^2)G'(z)=(z^2-8z-1)G(z).
\]
Comparison of coefficients yields, for \(n\ge0\),
\begin{equation}\label{rec}
 (n+3)R_{n+3}-(7n+13)R_{n+2}
       +(7n+15)R_{n+1}-(n+1)R_n=0.
\end{equation}

\paragraph{Explicit bounds.}
Put \(w=\sqrt2\), \(\rho=3+2w\), and
\begin{align*}
 P_\pm(x)&=64x^4-96x^3+(168-78w)x^2+(273w-537)x\pm640,\\
 L(x)&=\frac{P_-(x)}{64x^4},\qquad
 U(x)=\frac{P_+(x)}{64x^4}.
\end{align*}
Equivalently,
\[
 L(x),U(x)=1-\frac3{2x}
   +\frac{21/8-39w/32}{x^2}
   +\frac{-537/64+273w/64}{x^3}
   \mp\frac{10}{x^4}.
\]
For \(q_n=r_n/\rho\), recurrence \eqref{rec} becomes
\begin{equation}\label{T}
 q_{n+2}=T_n(q_n,q_{n+1}),\qquad
 T_x(a,b)=\frac{7x+13}{(x+3)\rho}
 -\frac{7x+15}{(x+3)\rho^2b}
 +\frac{x+1}{(x+3)\rho^3ab}.
\end{equation}

\begin{lemma}\label{bounds}
For every integer \(n\ge20\), \(L(n)<q_n<U(n)\).
Moreover, for every real \(x\ge20\),
\begin{align}
 &L(x)>0,\qquad (7x+15)\rho L(x)>x+1,\label{pos}\\
 &T_x(U(x),L(x+1))>L(x+2),\label{lower}\\
 &T_x(L(x),U(x+1))<U(x+2),\label{upper}\\
 &L(x)^2>U(x-1)U(x+1).\label{lc}
\end{align}
\end{lemma}

\begin{proof}
We first verify the four real-variable assertions by exact polynomial
arithmetic. The following five expressions are polynomials in
\(\mathbb Z[w][x]\):
\begin{align*}
 H_P(x)&=P_-(x),\\
 H_M(x)&=(7x+15)\rho P_-(x)-64(x+1)x^4,\\
 H_L(x)&=64(x+3)\rho^3(x+2)^4P_+(x)P_-(x+1)\\
 &\hspace{15mm}\cdot\bigl[T_x(U(x),L(x+1))-L(x+2)\bigr],\\
 H_U(x)&=64(x+3)\rho^3(x+2)^4P_-(x)P_+(x+1)\\
 &\hspace{15mm}\cdot\bigl[U(x+2)-T_x(L(x),U(x+1))\bigr],\\
 H_C(x)&=4096x^8(x-1)^4(x+1)^4
                  \bigl[L(x)^2-U(x-1)U(x+1)\bigr].
\end{align*}
Cancellation of the displayed denominators proves polynomiality.
Write
\[
 H_J(20+y)=\sum_{k=0}^{d_J}(a_{J,k}+b_{J,k}w)y^k.
\]
All \(a_{J,k},b_{J,k}\) are integers. Expansion using \(w^2=2\) gives
the following coefficient certificate:
\[
\begin{array}{c|r|r}
 J&d_J&
 \displaystyle\min_{0\le k\le d_J}
 \min\{a_{J,k}+\tfrac75b_{J,k},
       a_{J,k}+\tfrac32b_{J,k}\}\\ \hline
 P&4&64\\
 M&5&12672/5\\
 L&9&1574064128/5\\
 U&9&2525868032/5\\
 C&13&12288
\end{array}
\]
Since \(7/5<w<3/2\), every coefficient in each shifted polynomial is
strictly positive. Hence all five \(H_J(x)\) are positive for \(x\ge20\).
The first two give \eqref{pos}. Since \(P_+=P_-+1280\), the multipliers
defining \(H_L,H_U,H_C\) are positive on this range. The other three
give \eqref{lower}--\eqref{lc}.

For clarity, the coefficient table is a finite exact certificate,
not an asymptotic estimate. Its complete integer coefficient pairs
are recorded in \texttt{exact-certificate.json}.
Two independent implementations check them:
\texttt{verify\_proof.py} uses quadratic-field polynomial arithmetic,
and \texttt{verify\_independent.py} expands the products using only
integer and rational pairs and polynomial convolution.
The latter needs only Python's standard library.

It remains to propagate the bounds along the actual sequence.
Direct evaluation of the defining binomial sums gives
\[
 L(j)<q_j<U(j)\qquad(j=20,21).
\]
The same certificate records all four positive gaps.
For \(a\ge L(n)>0\) and \(b>0\), differentiation in \eqref{T} shows
\[
 \frac{\partial T_n}{\partial a}<0,\qquad
 \frac{\partial T_n}{\partial b}
 =\frac{(7n+15)\rho a-(n+1)}
       {(n+3)\rho^3ab^2}>0,
\]
where the second sign follows from \eqref{pos}.
If the bounds hold at \(n,n+1\), these signs imply
\[
 T_n(U(n),L(n+1))\le q_{n+2}
                  \le T_n(L(n),U(n+1)).
\]
Equations \eqref{lower} and \eqref{upper} therefore give the strict
bounds at \(n+2\). Induction starting with \(20,21\) completes the proof.
\end{proof}

\begin{proof}[Proof of Theorem \ref{main}]
For \(n\ge21\), all three indices \(n-1,n,n+1\) lie in the range of
Lemma \ref{bounds}. Thus
\[
 q_n^2>L(n)^2>U(n-1)U(n+1)>q_{n-1}q_{n+1}.
\]
Multiplication by \(\rho^2\) proves the desired inequality.
For \(5\le n\le20\), direct integer evaluation gives
\[
 \min_{5\le n\le20}
  \bigl(R_{n+1}^3R_{n-1}-R_n^3R_{n+2}\bigr)=4658935984>0.
\]
Every one of these sixteen exact integers is recorded in the
certificate. They are computed independently from the binomial sum
and from \eqref{rec}, starting with \(R_0=-1,R_1=1,R_2=7\).
This covers all \(n\ge5\).
Finally \(R_3=25,R_4=87,R_5=329,R_6=1359\), and
\[
 R_5^3R_3-R_4^3R_6=-4623352<0.
\]
Hence the ratio sequence cannot be extended back from index \(4\)
to index \(3\) while preserving log-concavity.
\end{proof}

\paragraph{Status and verification.}
The published full text retains Conjecture 4.7.
Current MathDB discussion and solution fields, focused literature
searches, and public claim catalogues were checked; no matching
solution claim was located. The separate literature audit records
the sources, timestamps, scope correction, and search limits.
The proof above covers the complete published statement; the
computations certify only its explicitly stated finite algebraic
obligations.

\begin{thebibliography}{9}
\bibitem{sun}
Brian Yi Sun,
\emph{On a ratio monotonicity conjecture of a new kind of numbers},
Journal of Inequalities and Applications \textbf{2018}, article 24.
\href{https://doi.org/10.1186/s13660-018-1614-1}{doi:10.1186/s13660-018-1614-1}.
Conjecture 4.7.
Earlier version: \href{https://arxiv.org/abs/1512.01008}{arXiv:1512.01008}.
\end{thebibliography}

\newpage
\paragraph{Task 2: determine the onset for the parameter family.}
For \(t>0\), define
\[
 R_n(t)=\sum_{k=0}^n\binom nk\binom{n+k}k\frac{t^k}{2k-1}.
\]
Determine the least \(N(t)\) for which the terms and consecutive
ratios are positive and the sequence
\[
 \{R_{n+1}(t)/R_n(t)\}_{n\ge N(t)}
\]
is log-concave.
Locate the parameter values at which \(N(t)\) changes.

\paragraph{Task 3: find sharp constants in the ratio bounds.}
With
\[
 B(n)=1-\frac3{2n}
   +\frac{21/8-39\sqrt2/32}{n^2}
   +\frac{-537/64+273\sqrt2/64}{n^3},
\]
determine the infimum and supremum of
\(n^4(q_n-B(n))\) over integers \(n\ge20\).
Prove where the extrema occur, including any limiting extremum.

\paragraph{Task 4: quantify higher orders of log-monotonicity.}
For each integer \(k\ge2\), determine the least threshold \(N_k\)
such that
\[
 (-1)^k\Delta^k\log R_n>0\qquad(n\ge N_k),
 \qquad \Delta f(n)=f(n+1)-f(n).
\]
Obtain effective bounds for \(N_k\) as \(k\) grows.

\paragraph{Task 5: find a structural proof of the ratio inequality.}
Develop an integral representation, a continued fraction, or a
positivity-preserving operation that explains the strict ratio
log-concavity directly. Determine which other sequences with
generating functions
\(-\sqrt{1-az+z^2}/(1-z)^b\) admit the same argument.
\end{document}
